\section{Finite Gaussian quadrature}
\label{sec:huang_gaussian_quadrature}

The Gaussian error-reduction construction of
\cite[Section 5.4, Lemma 19]{Huang2015GroundState} separates truncation in time,
finite quadrature and approximation of each propagator. The results below establish
the quadrature step with exact Gaussian interval masses. They do not construct a
matrix product operator approximation to the full time evolution.

\begin{definition}[Time mesh and Gaussian weights]
    \label{def:huang_time_mesh}
    \lean{TNLean.MPS.AreaLaw.timeMeshCount,TNLean.MPS.AreaLaw.timeMeshCell,
        TNLean.MPS.AreaLaw.timeMeshNode,TNLean.MPS.AreaLaw.timeMeshWeight}
    \leanok
    Let $T\ge0$, $\delta>0$, and let $\gamma_h$ be the centered Gaussian measure
    of variance $h\ge0$. Put $N=\lceil2T/\delta\rceil+1$ and, for $0\le j<N$, define
    \begin{align*}
        I_j & =[-T,T]\cap[-T+j\delta,-T+(j+1)\delta),                         \\
        t_j & =\min\{-T+j\delta,T\},                  & w_j & =\gamma_h(I_j).
    \end{align*}
    At $h=0$, the Gaussian is the unit mass at zero.
\end{definition}

\begin{lemma}[Mass and number of samples]
    \label{thm:huang_time_mesh_mass}
    \lean{TNLean.MPS.AreaLaw.timeMeshNode_mem,TNLean.MPS.AreaLaw.timeMeshCount_le,
        TNLean.MPS.AreaLaw.exists_timeMeshNode_abs_sub_le,
        TNLean.MPS.AreaLaw.timeMeshWeight_nonneg,TNLean.MPS.AreaLaw.sum_timeMeshWeight,
        TNLean.MPS.AreaLaw.sum_timeMeshWeight_le_one}
    \leanok
    \uses{def:huang_time_mesh}
    Every node belongs to $[-T,T]$, and every point of that interval lies
    within $\delta$ of a node. The weights satisfy
    \[
        w_j\ge0,\qquad \sum_jw_j=\gamma_h([-T,T])\le1,
        \qquad N\le2T/\delta+2.
    \]
\end{lemma}
\begin{proof}
    \leanok
    The cells are pairwise disjoint and cover $[-T,T]$, including the right endpoint.
    Choosing a cell containing a given point gives its nearby node.
    Additivity gives the mass identity. The node bounds follow from their definition,
    and the ceiling estimate gives the sample count.
\end{proof}

\begin{lemma}[Gaussian quadrature error]
    \label{thm:huang_time_mesh_error}
    \lean{TNLean.MPS.AreaLaw.norm_setIntegral_sub_timeMesh_sum_le,
        TNLean.MPS.AreaLaw.norm_integral_sub_timeMesh_sum_le}
    \leanok
    \uses{def:huang_time_mesh}
    Let $f$ take values in a real Banach space, be integrable for $\gamma_h$, and satisfy
    $\|f(t)-f(s)\|\le M|t-s|$ for $s,t\in[-T,T]$, with $M\ge0$. Then
    \[
        \left\|\int_{-T}^T f(t)\,d\gamma_h(t)-\sum_jw_jf(t_j)\right\|
        \le M\delta.
    \]
    If $\|f(t)\|\le1$ for every $t$, then
    \[
        \left\|\int_{\mathbb R}f(t)\,d\gamma_h(t)-\sum_jw_jf(t_j)\right\|
        \le2e^{-T^2/(2h)}+M\delta.
    \]
    For the last estimate at $h=0$, the exponential term is understood as $2$.
\end{lemma}
\begin{proof}
    \leanok
    \uses{thm:huang_time_mesh_mass}
    For $t\in I_j$ one has $|t-t_j|\le\delta$. Integrate the Lipschitz estimate
    on each cell and sum; the total Gaussian mass is at most one. The full integral
    differs from the truncated integral by at most the Gaussian mass outside $[-T,T]$,
    which is bounded by the displayed exponential term.
\end{proof}

\begin{lemma}[Quadrature at prescribed mesh error]
    \label{thm:huang_time_mesh_accuracy}
    \lean{TNLean.MPS.AreaLaw.timeMeshCount_accuracy_le,
        TNLean.MPS.AreaLaw.norm_integral_sub_timeMesh_accuracy_sum_le}
    \leanok
    \uses{thm:huang_time_mesh_error}
    Let $\varepsilon>0$. Taking $\delta=\varepsilon/(M+1)$ gives error at most
    $2e^{-T^2/(2h)}+\varepsilon$, with
    \[
        N\le\frac{2T(M+1)}{\varepsilon}+2.
    \]
    The same nodes and weights work for every admissible integrand.
\end{lemma}
\begin{proof}
    \leanok
    Since $M\delta\le\varepsilon$, the preceding theorem applies. Substitute the
    mesh into the sample-count bound.
\end{proof}

\begin{lemma}[Quadrature of Hermitian time evolution]
    \label{thm:huang_unitary_quadrature}
    \lean{TNLean.MPS.AreaLaw.norm_hermitianUnitaryPath_sub_time_le,
        TNLean.MPS.AreaLaw.norm_gaussianIntertwiner_sub_unitary_timeMesh_sum_le}
    \leanok
    \uses{thm:huang_time_mesh_accuracy}
    Let $H$ be Hermitian on a nonzero finite-dimensional Hilbert space. Its unitary path
    $U(t)=e^{itH}$ satisfies $\|U(t)-U(s)\|\le\|H\||t-s|$. Consequently
    \[
        \left\|\int_{\mathbb R} e^{itH}\,d\gamma_h(t)-\sum_jw_je^{it_jH}\right\|
        \le2e^{-T^2/(2h)}+\varepsilon
    \]
    for the mesh $\varepsilon/(\|H\|+1)$, with at most
    $2T(\|H\|+1)/\varepsilon+2$ samples. This estimate applies in particular to a
    Hamiltonian shifted by an energy estimate.
\end{lemma}
\begin{proof}
    \leanok
    The group law for the unitary path gives
    $\|U(t)-U(s)\|=\|U(t-s)-I\|$. The Duhamel estimate against the zero generator
    bounds this by $\|H\||t-s|$. Unitarity bounds the integrand norm by one, so the
    preceding theorem applies in operator norm.
\end{proof}
